\track{T8 v1.1 修补：盲测点名的五处「推不动」}{按两轮盲测反馈定向补写：每条给出显式步骤、定理清单或约定说明，不重复 T1--T7。}

\begin{word}{从 S-对偶到二维镜像：三步}{S-duality to two-dimensional mirror symmetry}{sduality-to-mirror}{\S2.2, p.~11--15; \S3.5, p.~28--31; \S4.4, p.~52--55; \S5.2, p.~65--68; \S5.4--5.5, p.~72--77}
\mathdef{取 $g>1$ 的 $C$ 与单链单连通 $G$。四维 $\mathcal N=4$ 的 $\tau=\theta/2\pi+4\pi\mathrm{i}/e^2$；对偶群由 $S:(\tau,G)\mapsto(-1/(n_g\tau),{}^LG)$、$T:\tau\mapsto\tau+1$ 生成（§2.2, p.~11）。二维真正依赖的是典范参数 $\Psi=\operatorname{Re}\tau+\mathrm{i}\operatorname{Im}\tau\frac{t-t^{-1}}{t+t^{-1}}$（(3.53)），$t\in\mathbb{CP}^1$ 为 GL 扭化参数；$S$ 给 $\Psi\mapsto-1/(n_g\Psi)$（(5.35)）。}
\physdef{$S$ 在四维交换 Wilson 与 't Hooft 算子、交换强弱耦合。紧化在 $C$ 上，低能理论以 $\mathcal M_H(G,C)$ 为靶；二维只用到 $\operatorname{Im}\tau$ 的整体缩放与 $\theta$ 的拓扑项 $B=-\operatorname{Re}\tau\,\omega_I$，二者合成 (5.34)：$\tau$ 不是好参数，$\Psi$ 才是。}
\mathex{三步。一：$S$ 把 $(G,\tau)$ 送到 $({}^LG,-1/(n_g\tau))$，$\Psi$ 到 $-1/(n_g\Psi)$，$0\leftrightarrow\infty$；$T$ 给 $\Psi\mapsto\Psi+1$（(5.36)）。二：Hitchin 纤维化 $\pi:\mathcal M_H\to\mathcal B$ 的纤维关于 $\omega_J,\omega_K$ 都 Lagrangian，复维数 $\dim\mathcal B=(g-1)\dim G$（§4.4, (4.25)），是复环面；$S$ 在基上给 $\Xi:\mathcal B\to{}^L\mathcal B$。三：$\mathcal M_H(G)$ 与 $\mathcal M_H({}^LG)$ 的一般纤维互为对偶复环面，沿纤维的 $T$-对偶把 Kähler 模与复模互换；对超凯勒流形这正是 $J$ 与 $K$ 对调、$I$ 不变，于是 $G$ 侧的 $\omega_K$ 在对偶侧成为复结构数据（$\omega_J$ 与 $\omega_K$ 互换），故依赖 $\omega_K$ 的 A-模型变依赖 $J$ 的 B-模型。}
\physex{特殊值：$t=\pm\mathrm{i}$ 给 $\Psi=\infty$，即 $J$ 中 B-模型；$t=\pm1$ 给 $\Psi=\operatorname{Re}\tau$，即 $K$ 中 A-模型。$S$ 把 $t$ 映为 $-t(\tau/|\tau|)$，$\tau=\mathrm{i}$ 时 $\mathrm{i}\mapsto1$（§5.5, p.~74）。原文选的就是这一对：$G$ 侧 $t=1,\theta=0$ 得 $(\mathcal M_H(G),\omega_K)$ 上 A-模型，${}^LG$ 侧取 ${}^Lt=\mathrm{i},{}^{L}\theta=0$ 得 $(\mathcal M_H({}^LG),J)$ 上 B-模型。}
\bridge{对应规则：$\tau\leftrightarrow$ 四维耦合，$\Psi\leftrightarrow$ 二维典范参数，$S\leftrightarrow\Xi$ 加纤维 $T$-对偶，$\omega_K$ 的 A-模型 $\leftrightarrow$ $J$ 的 B-模型。$\Xi$ 一般非恒等：$v$ 由 $H_\alpha-h_\alpha=0$ 定义时 $\Xi(v)$ 由 ${}^LH_\alpha-h_\alpha=0$ 定义，单链群上恰为恒等，故 SYZ 中默认平凡的 $\Xi$ 在此确可平凡。原文只断言「基上经典作用、纤维上经 $T$-对偶作用」（§5.4, p.~72），纤维对偶性归于 SYZ，细节据 Witten 2008 与本词典整理；§5.5 末（p.~77）明说未处理 $G$ 的中心与 $C$ 的自旋结构。}
\intuit{一句话：四维 $S$ 换群又换耦合；到二维，换群成基上重排，换耦合成纤维上 T-对偶，而 T-对偶把辛结构换成复结构，A 与 B 遂对调。}
\pitfall{三处易错。一：纤维不是 $T^{2r}$，其复维数是 $(g-1)\dim G$（$SU(2)$ 时 $3g-3$），$r$ 只经 $\dim G$ 进场。二：$\Psi$ 非新输入，只是 $\tau$ 与 $t$ 的组合；$t=\pm\mathrm{i}$ 处公式有极点，须改用 $J$ 中 B-模型。三：「$S$ 即镜像」非一步：$S$ 只给「$J$ 中 B $\cong$ $K$ 中 A」，要两边同在 $J$ 还须复合 $\mathcal I:\varphi\mapsto\mathrm{i}\varphi$，它把理论移出 $\mathbb{CP}^1_g$ 族。细则见 T7『三种拓扑扭化』条。}
\end{word}

\begin{word}{左右约定总表}{a table of left--right conventions}{conventions-g-lg}{\S3.1, p.~17--21; \S5.2--5.6, p.~65--77; \S6.1, p.~79; \S6.4, p.~95--102; \S8, p.~116--122; \S9.1--9.3, p.~132--148}
\mathdef{记 $G$ 为规范群，${}^LG$ 为其对偶。$G$ 侧是 $\operatorname{Bun}_G$ 与 $\mathcal M(G,C)$；${}^LG$ 侧是平坦 ${}^LG_{\mathbb C}$-丛的模空间 $\operatorname{Loc}_{{}^LG}$，即 $\mathcal M_H({}^LG,C)$ 的复结构 $J$。几何 Langlands：$\operatorname{Loc}_{{}^LG}$ 上相干层与 $\operatorname{Bun}_G$ 上扭化 $\mathcal D$-模对偶，点 $E$ 的天穹层 $\mathcal O_E$ 对应 Hecke 特征层 $\mathcal F_E$（§6.4, p.~95）。}
\physdef{(a) 电与磁。正确：电特征膜 $E_p(R)$ 是 $\Psi=\infty$ 处 Wilson 算子的共同特征膜，特征值取 $G$ 的表示 $R$；磁特征膜是 $\Psi=0$ 处 't Hooft 算子的共同特征膜，特征值取 ${}^LG$ 的 ${}^LR$。$S:\Psi\mapsto-1/(n_g\Psi)$ 交换 $\infty$ 与 $0$，故交换电与磁、也交换 $G$ 与 ${}^LG$（§6.4, p.~97）。易反：把电配给 ${}^LG$；检验看 OPE——Wilson 系数来自 $G$ 侧张量积，'t Hooft 来自 ${}^LG$ 侧。}
\mathex{(b) 零膜与层。正确：$x\in\mathcal M_H(G,C)$ 上零膜的 Chan-Paton 丛是天穹层 $\mathcal O_x$；取 $J$ 时 $x$ 是平坦 $G_{\mathbb C}$-丛，故零膜自带 $\operatorname{Loc}_G$ 侧数据，这正是「电特征膜即零膜」（§8, p.~116）。换成 ${}^LG$ 侧的平坦丛 $E$ 后，天穹层 $\mathcal O_E$ 就是那边的电特征膜；对它作 $S$-对偶得 $G$ 侧缠 Hitchin 纤维的 type $F$ 磁特征膜，再经翻译器成 $\operatorname{Bun}_G$ 上 $\mathcal F_E$。易反：把「$E$ 处结构层 $\mathcal O_E$」直接读成 $\operatorname{Bun}_G$ 上的层；$\mathcal O_E$ 住 $\operatorname{Loc}_{{}^LG}$，$\mathcal F_E$ 住 $\operatorname{Bun}_G$，中间隔着 $S$-对偶与翻译两步。}
\physex{(c) 两类膜互换。按 $(I,J,K)$ 顺序：$S$ 把 $(B,B,B)$ 送到 $(B,A,A)$、$(B,A,A)$ 送回 $(B,B,B)$、$(A,B,A)$ 送到自己；逐因子 $J$ 中 B-膜$\mapsto$ $K$ 中 A-膜（§5.2, Table 2；§5.6, p.~77）。故「$J$ 中 B」与「$K$ 中 A」互为对偶；c.c. 膜是 $(A,B,A)$，$S$ 保其类型（§12.4）。易反：说成「$J$ 中 B 对偶于 $J$ 中 A」——那是再复合 $\mathcal I$ 的镜像说法。}
\bridge{(d) 记号。原文 §11.1（p.~185）记 $\mathcal B_{c.c.}$，条件 $F=(\operatorname{Im}\tau)\omega_J$、$(\omega^{-1}F)^2=-1$；文献亦作 $A_{cc}$，其中 $A$ 指「$K$ 扭化中的 A-膜」。同一对象：$\mathcal B$ 只是 brane 的首字母，与 A/B-模型无关。本词典统一写 $\mathcal B_{cc}$。易反：把 $\mathcal B_{cc}$ 的 B 读成「B-模型」，得出「它是 B-膜」；实际它在 $K$ 中是 A-膜、在 $J$ 中才是 B-膜。}
\intuit{一句话：$G$ 侧（$\operatorname{Bun}_G$）看到 D-模与特征层，${}^LG$ 侧（$\operatorname{Loc}_{{}^LG}$）看到零膜与结构层；$S$ 负责换边，翻译器负责把膜换成层，两步不可并成一句。}
\pitfall{盲测点名的张力：5.8 条把电特征膜写成 $\operatorname{Bun}_G$ 上的 Hecke 特征层，6.1 条写成 $\operatorname{Loc}$ 侧结构层 $\mathcal O_E$、其镜像 A-膜才译成 $\mathcal F_E$。两句都对但主语不同：前者说它在 $G$ 理论里的像，后者说它本身。凡遇「电特征膜」先问在哪个理论里，凡遇「特征层」先问在 $\operatorname{Bun}$ 还是 $\operatorname{Loc}$。口诀：电--Wilson--$G$--$\Psi=\infty$ 对磁--'t Hooft--${}^LG$--$\Psi=0$。细则见 T7『膜到 D-模的翻译器』条。}
\end{word}

\begin{word}{翻译器的构造、命题与证明状态}{the translator: construction, statement, status}{translator-status}{\S11.1, p.~185--189; \S11.2, p.~189--200; \S11.3, p.~200--206}
\mathdef{对象是 $K$ 扭化下 $\mathcal M_H(G,C)$ 的 A-膜范畴。构造：空间充满的 A-膜 $\mathcal B_{cc}$，Chan-Paton 为线丛 $\mathcal N$ 带酉联络 $\alpha$、$F=\mathrm{d}\alpha$，A-膜条件 $(\omega^{-1}F)^2=-1$（(11.1)），$\omega=(\operatorname{Im}\tau)\omega_K$；解 $F=(\operatorname{Im}\tau)\omega_J$ 给出 $\hat N=\omega^{-1}F=\omega_K^{-1}\omega_J=I$，且 $\omega-\mathrm{i}F$ 为 $(2,0)$ 型全纯辛形式（(11.2)）。命题：$\mathcal A\mapsto\operatorname{Hom}(\mathcal B_{cc},\mathcal A)$ 给出该范畴到 $\operatorname{Bun}_G$ 上扭化 $\mathcal D$-模的等价。}
\physdef{构造的四件数据须凑齐。线丛整体存在，因 $\omega_J=\mathrm{d}\alpha$ 且周期为 $2\pi$ 整数倍，显式 $\alpha=\frac{\operatorname{Im}\tau}{2\pi}\int_C|\mathrm{d}^2z|\operatorname{Tr}(\phi_z\delta A_{\bar z}+\phi_{\bar z}\delta A_z)$（(11.3)）；同族 $(\operatorname{Im}\tau)(a\omega_J+b\omega_I)$、$a^2+b^2=1$，$b\neq0$ 时周期非整、不是任何联络的曲率。开弦代数：圆盘瞬子被解析性逼成常数，故无瞬子修正，鬼数零的边界算子即关于 $\hat N=I$ 的全纯函数层，$F\neq0$ 时为按 $\Omega_I$ 的 Poisson 括号的非交换形变。}
\mathex{两步代数化。作用量化为边界项 $\mathcal S=-\mathrm{i}\operatorname{Im}\tau\int_{\partial\Sigma}\Phi^*\varpi$（(11.8)--(11.10)），弦代数为 $\mathcal D_{K^{1/2}}$。为何是 $K^{1/2}$：时间反演把代数换成相反代数，而 $\mathcal D_L^{\mathrm{op}}\cong\mathcal D_{L^{-1}\otimes K_M}$，自相反要求 $L^2\cong K_M$，故 $L=K_M^{1/2}$——这是纯代数条件；换平方根即换平凡化，差别是 $\mathbb{Z}_2$-gerbe 扭化。}
\physex{按 $\Psi$ 细化：一般 $\Psi$ 有 $\mathcal B_{cc}^\Psi$，开弦代数 $\mathcal D_{K^{1/2}\otimes\mathcal L^\Psi}$（$\mathcal L$ 为行列式线丛）；一般 $\Psi$ 处 Hitchin 纤维不是 A-膜（$F+B|_Y=0$ 无解，(11.29)）。$\Psi=h$（对偶 Coxeter 数）时退化为普通 $\mathcal D$-模；把 $S$ 与 $T^h$ 复合，就把 $\mathcal M_H({}^LG)$ 的零膜映成 $\mathcal M(G,C)$ 上未扭化 $\mathcal D$-模，这才是几何 Langlands 的标准表述。}
\bridge{证明状态逐条。构造与无瞬子修正：§11.3 的动词是 argue；开弦代数 $=\mathcal D_{K^{1/2}}$ 是物理推导，数学对应 Nadler--Zaslow；「A-膜范畴 $\cong$ 非交换形变上 B-膜范畴」引 Kapustin，仍 argued；到 $\operatorname{Bun}_G$ 叠的严格化原文只说必须，未做。已知定理仅特例：$GL_1$ 为 Fourier--Mukai（Laumon--Rothstein），$GL_n$ 由 Frenkel--Gaitsgory 型工作，oper 由 Beilinson--Drinfeld；一般群的等价（Frenkel (3.1)）仍是猜想。}
\intuit{一句话：$\mathcal B_{cc}$ 是全空间的真空，其他 A-膜是它上面的波函数；配上到真空的开弦态，就翻成 D-模。}
\pitfall{三处易错。一：$\operatorname{Hom}(\mathcal B_{cc},\mathcal A)$ 是复形或层，取上同调才得向量空间，等价只在导出范畴层面成立。二：$K_M^{1/2}$ 是模空间 $\mathcal M$ 上典范丛的平方根，与 $C$ 上的自旋结构是两回事；它由 $L^2\cong K_M$ 逼出，换平方根即换 $\mathbb{Z}_2$-gerbe 平凡化。三：以为 §11 证明了等价：原文是 argued，且 $\mathcal B_{cc}$ 只在 $\mathcal M_H$ 光滑部分定义、层化只在 $T^*\mathcal V$ 型开集可行。细则见 T7『膜到 D-模的翻译器』条。}
\end{word}

\begin{word}{单极子模空间即 Hecke 修正空间}{monopole moduli as Hecke modification spaces}{monopole-hecke-theorem}{\S9.1--9.2, p.~132--148; \S10.1--10.2, p.~148--155; \S10.4, p.~169--180; \S10.5, p.~182--183}
\mathdef{$p_0\in C$，${}^Lw$ 为 ${}^LG$ 支配余权，$\mathcal Y({}^Lw,p_0)$ 为 $p_0$ 处型 ${}^Lw$ 的 Hecke 修正空间（丛加 $p_0$ 外同构），$\bar{\mathcal Y}$ 是紧化（§9.2, p.~145）。$\mathcal M^{\mathrm{mon}}$ 为 $W=I\times C$ 上带 't Hooft 奇点与椭圆边条件的 Bogomolny 解模空间。命题：$\mathcal M^{\mathrm{mon}}\cong\bar{\mathcal Y}({}^Lw;p_0)$，且 $\bar{\mathcal Y}$ 是仿射 Grassmann 的 Schubert 簇（§10.5, p.~182）。}
\physdef{边界条件（§10.1, p.~148）：$E_-$ 端 $A_-$ 取 Dirichlet、$\phi_0$ 自由，$E_+$ 端 $\phi_0=0$、$A$ 取协变 Neumann——每端在定 $A$ 与定 $\phi_0$ 中二选一才椭圆。解取 $y=$ 常数切片就是一次修正；反之由 $\partial_yA_{\bar z}=-\mathrm{i}D_{\bar z}\phi_0$ 沿 $y$ 唯一延拓成解。故等同是构造性的。}
\mathex{定理四层。一，Bogomolny 解与 Nahm 数据：ADHMN 给公式（Atiyah--Drinfeld--Hitchin--Manin 1978），Nahm（1979）提出，Hitchin（1983）对 $SU(2)$ 证一一对应。二，与单极子模空间的超凯勒等距：Donaldson（1984）对 $SU(2)$，Hurtubise 推广。三，同构于 Schubert 簇：§10.5 引用，即仿射 Grassmann 的轨道闭包（Beilinson--Drinfeld 框架）；§9.3 指出 Schubert 胞腔即 $\mathcal Y(\rho)$。四，OPE 系数由相交上同调给出：§10.4 只是物理推导，把 Lusztig--Ginzburg--Mirkovi\'c--Vilonen 的几何 Satake 当外部定理；原文按语指出难点是证 $\tilde N_\alpha\cong N_\alpha$。}
\physex{例子。$G=U(2)$、${}^Lw=(1,0)$：修正只依赖 $p_0$ 处截面的复缩放，$\mathcal Y\cong\mathbb{CP}^1$，紧而光滑（§9.2）。${}^Lw=(2,0)$：未紧化是挖一点的 $\mathbb{P}(1,1,2)$，紧化补入型 $(1,1)$ 的修正，得 $\bar{\mathcal Y}\cong\mathbb{P}(1,1,2)$，奇点局部为 $\mathbb{C}^2/\mathbb{Z}_2$。荷 $k$ 的 $SU(2)$ 单极子模空间实维数 $4k$（$k=1$ 为 $\mathbb{R}^3\times S^1$）；指标公式给单个算子的复维数贡献 $a$（Pauly 定理；原文脚注：其 $k$ 即本文 $a/2$）。}
\bridge{修正空间复维数 $a$，单极子模空间实维数 $4k$；由 $a=2k$，$k=1$ 对应 ${}^Lw=(2,0)$，复 2 维与 $4k=4$ 相符。$(1,0)$ 是 $a=1$ 即 $k=1/2$ 的半情形，拿它对 $4k$ 会差因子 2，这不是矛盾，是权与荷的换算。要把 $4k$ 维的模空间降至 §10，还须固定质心 $\mathbb{R}^3$（含 $y$ 方向）与整体 $S^1$；此步原文未对账，属本词典整理。}
\intuit{一句话：'t Hooft 算子把磁荷钉在一点；允许的奇点形状就是单极子，一个解的 $y=$ 常数截面与一次修正本是同一份数据。}
\pitfall{三处易错。一：把等同当定理——等同是定义式的，真正的定理只是奇点外光滑、维数等于椭圆指标与 Schubert 簇识别。二：拿 $(1,0)$ 的 $\mathbb{CP}^1$ 对 $4k$ 得「维数不符」，漏了 $a=2k$。三：把 Nahm 数据当原文正文——§2.1 的「Nahm 定理」是十维超对称的维数定理，与单极子的 Nahm 方程无关。细则见 T7『单极子模空间与 Hecke 修正的机制』条与『Nahm 极点与 ADHMN 构造』条。}
\end{word}

\begin{word}{§12 四维提升的实质内容}{branes from gauge theory: what §12 actually does}{sigma12-lift}{\S12.1, p.~208--211; \S12.2, p.~211--215; \S12.3, p.~215--218; \S12.4, p.~218--222}
\mathdef{取带边流形 $M=\mathbb{R}^2_+\times C$，$x^1=0$ 为边界，膜即边界条件。判据是超流法向分量为零：$\bar\epsilon J_1=0$（(12.7)）。先定费米子条件 $T\lambda|=\pm\lambda|$，由 $\delta\lambda=\frac12\Gamma^{IJ}F_{IJ}\epsilon$（(12.8)）推玻色条件，再回验 $J_1=0$。记账：四维 16 超荷保 4、二维 8 超荷保 4，故四维 $1/4$ BPS 即二维 $1/2$ BPS。}
\physdef{三类边条件（§12.2--12.4）。$(B,B,B)$：$T=\Gamma_0\Gamma_4$ 逼出场强自对偶，即 Hitchin 方程 (12.16)，边界值是 $\mathcal M_H$ 的点，即零膜。$(B,A,A)$：对 $(B,B,B)$ 施加 $S$，投影算子给 $T'=\frac{1+\hat\Gamma}{\sqrt2}T\frac{1-\hat\Gamma}{\sqrt2}=\Gamma_1\Gamma_2\Gamma_3\Gamma_4$（(12.21)--(12.23)），特征多项式固定而 $A_{\bar z}$ 可涨落，膜缠 Hitchin 纤维。$(A,B,A)$：$T'=-\Gamma_{15}\Gamma_{2-7}\Gamma_{3+6}$（(12.28)），固定 $\mathcal A_{\bar z}=A_{\bar z}+\mathrm{i}\phi_{\bar z}$ 的子流形 $W'$，在 $J$ 中复、关于 $\Omega_J$ Lagrangian。}
\mathex{$S$ 的像：$(B,B,B)\leftrightarrow(B,A,A)$，原文说「$S$ 把缠 Hitchin 纤维的 $(B,A,A)$ 映为零膜」，阿贝尔情形可验；$(A,B,A)$ 自对偶——对 (12.28) 施 $S$ 得 $-\frac{\Gamma_{15}}{2}((\Gamma_{26}+\Gamma_{37})-\hat\Gamma(\Gamma_{23}+\Gamma_{67}))$（(12.30)），算后与混合 Dirichlet--Neumann 条件 (12.33) 比较，读出 $f=\omega_J$（(12.34)），正是 c.c. 膜。三类膜的像由 Table 2 决定。}
\physex{只覆盖非分歧与光滑部分的原因。$C$ 上无插入点即非分歧，带奇点者另在 §4 与 §11.2。可约点边界值不唯一：$(B,B,B)$ 另有 $T=\Gamma_{023467}$ 的构造（(12.17)--(12.19)），只保 4 超荷，光滑点等价、可约点不等价。$(B,A,A)$ 的经典分析给 $[\phi_2,\phi_3]=0$，与完整方程差 $F_{23}$，原文猜后者是 RG 流修正。时间反演使 CP 丛只能平凡，原文自认构造不完整。}
\bridge{与数学定理的距离：§12 是纲要——它把膜装回四维、给出三类边条件与 $\Gamma$ 矩阵，但未证明范畴等价。可对照的严格结果只是特例：非分歧 $GL_n$ 由 Frenkel--Gaitsgory 型工作，oper 由 Beilinson--Drinfeld 用临界层仿射 Kac--Moody 构造，而 oper 是谱参数取特殊值的一族；一般群仍是猜想。KW06 的贡献是把「特征 D-模」重释为四维边界条件在二维有效理论中的像，给出构造方案而非证明。}
\intuit{一句话：§12 不靠玩具模型——膜是四维规范理论的边界条件，由四维保证。}
\pitfall{三处易错。一：以为 §12 证明了范畴等价——原文自述纲要，且自点可约点、$[\phi_2,\phi_3]$ 与 CP 丛三处不完整。二：把 $D_{\bar z}\phi_z=0$ 当完整 Hitchin 方程——它只是 $J$ 中方程的全纯半支，正因如此膜落在纤维上而非空间充满。三：以为 $(A,B,A)$ 在 $S$ 下会换类——它自对偶，因为读出的是不变的 $f=\omega_J$。细则见 T7『三种拓扑扭化』条。}
\end{word}
